% ============ 18 紧支集上同调群 (原书 pp.118-122) ============
\chapter{紧支集上同调群}\label{sec:18}
\index{紧支集上同调}
\index{完备化}
\index{紧支集高阶直像}
\index{Nagata 嵌入定理}

\subsection*{来自拓扑的启发}

在平展上同调中，与在拓扑中一样，定义紧支集上同调群很重要---我们在曲线的情形已经看到这一点(\S 14).它们应当与通常的上同调群对偶.

传统定义(Greenberg 1967, p.162)是，对流形 $U$，
\[
H^{r}_{c}(U, \mathbb{Z}) = \varinjlim_{Z} H^{r}_{Z}(U, \mathbb{Z}),
\]
其中 $Z$ 遍历 $U$ 的紧子集.更一般地(Iversen 1986, III.1)，当 $F$ 是局部紧拓扑空间 $U$ 上的层时，定义
\[
\Gamma_c(U, F) = \varinjlim_{Z} \Gamma_Z(U, F),
\]
其中 $Z$ 同样遍历 $U$ 的紧子集，并设 $H^{r}_{c}(U, \cdot) = R^{r}\Gamma_c(U, \cdot)$.

对代数簇 $U$ 与 $U_{\mathrm{et}}$ 上的层 $F$，这提示定义
\[
\Gamma_c(U, F) = \varinjlim_{Z} \Gamma_Z(U, F),
\]
其中 $Z$ 遍历 $U$ 的完备子簇 $Z$，并设 $H^{r}_{c}(U, \cdot) = R^{r}\Gamma_c(U, \cdot)$.

然而，这个定义导致异常的群.例如，若 $U$ 是代数闭域上的仿射簇，则 $U$ 的完备子簇只有有限子簇(AG 7.5)，而对有限子簇 $Z \subset U$，
\[
H^{r}_{Z}(U, F) = \bigoplus_{z \in Z} H^{r}_{z}(U, F).
\]
因此，若 $U$ 光滑且维数为 $m$，$\Lambda$ 是常值层 $\mathbb{Z}/n\mathbb{Z}$，则
\[
H^{r}_{c}(U, \Lambda) = \varinjlim H^{r}_{Z}(U, \Lambda) = \bigoplus_{z \in U} H^{r}_{z}(U, \Lambda) = \begin{cases}
\bigoplus_{z \in U} \Lambda(-m), & r = 2m;\\
0, & \text{其他情形}.
\end{cases}
\]
这些群甚至不有限.我们需要不同的定义.

若 $j \colon U \to X$ 是拓扑空间 $U$ 到局部紧空间 $X$ 的开子集上的同胚，则
\[
H^{r}_{c}(U, F) = H^{r}(X, j_{!}F)
\]
(Iversen 1986, p.184).特别地，当 $X$ 紧时，
\[
H^{r}_{c}(U, F) = H^{r}(X, F).
\]
我们把这个作为我们的定义.

\subsection*{紧支集上同调}

\begin{definition}\label{def:18.1}
对簇 $U$ 上的任何挠层 $F$，定义
\[
H^{r}_{c}(U, F) = H^{r}(X, j_{!}F),
\]
其中 $X$ 是含 $U$ 为其稠密开子簇的任何完备簇，$j$ 是包含映射.
\end{definition}

开浸入 $j \colon U \hookrightarrow X$(从 $U$ 到完备簇 $X$，使 $j(U)$ 在 $X$ 中稠密)称为 $U$ 的完备化(或紧化).沿用拓扑中的术语，我们把 $H^{r}_{c}(U, F)$ 称为 $F$ 的紧支集上同调群(而不是更合逻辑的完备支集上同调群).

这个定义\footnote{对十九世纪迷信的信徒而言，他们认为所有簇都天然地带有射影空间中的嵌入，两个问题都不会出现.}提出两个问题：
\begin{enumerate}
  \item 每个簇都有完备化吗？
  \item 紧支集上同调群不依赖于完备化吗？
\end{enumerate}
第一个问题由 Nagata 在 1962 年给出肯定回答.\footnote{Nagata, Masayoshi, \emph{Imbedding of an abstract variety in a complete variety}, J. Math. Kyoto Univ. 2 (1962), 1--10.\\
Nagata, Masayoshi, \emph{A generalization of the imbedding problem of an abstract variety in a complete variety}, J. Math. Kyoto Univ. 3 (1963), 89--102.}更一般地，他证明：对从一个 Noether 概形到另一个的有限型分离态射 $\varphi \colon U \to S$，存在紧合态射 $\bar{\varphi} \colon X \to S$ 与开浸入 $j \colon U \to X$，使得 $\varphi = \bar{\varphi} \circ j$:
\[
\begin{tikzcd}
U \arrow[rr, "j"] \arrow[rd, "\varphi"'] & & X \arrow[ld, "\bar{\varphi}"] \\
& S &
\end{tikzcd}
\]
Nagata 的原始证明用赋值环.更现代，概形论的证明见 L\"utkebohmert, W., \emph{On compactification of schemes}, Manuscripta Math. 80 (1993), 95--111，以及 Deligne, Pierre, \emph{Le th\'eor\`eme de plongement de Nagata}. Kyoto J. Math. 50 (2010), no. 4, 661--670.

注意，与曲线的情形不同，在高维中嵌入不唯一：从任何完备化出发，都可以通过爆破与压平边界的子簇构造其他的完备化.然而，下一个命题表明 (b) 的答案也是肯定的---需要要求 $F$ 是挠层，以便能应用紧合基变换定理.

\begin{proposition}\label{prop:18.2}
当 $F$ 是挠层时，群 $H^{r}(X, j_{!}F)$ 不依赖于嵌入 $j \colon U \hookrightarrow X$ 的选取.
\end{proposition}

\begin{proof}
设 $j_1 \colon U \hookrightarrow X_1$ 与 $j_2 \colon U \hookrightarrow X_2$ 是 $U$ 的两个完备化.考虑到乘积的对角映射 $j \colon U \hookrightarrow X_1 \times X_2$，设 $X$ 是 $U$ 在 $X_1 \times X_2$ 中的闭包.则 $j \colon U \hookrightarrow X$ 是 $U$ 的完备化，而投影映射是 $X \to X_1$，$X \to X_2$，在 $U$ 上诱导恒等映射.只需证明 $H^{r}_{c}(X_1, j_{1!}F) \cong H^{r}_{c}(X, j_{!}F)$.考虑：
\[
\begin{tikzcd}
U \arrow[r, "j"] \arrow[rd, "j_1"'] & X \arrow[d, "\varphi"] \\
& X_1
\end{tikzcd}
\]
Leray 谱序列给出
\[
H^{r}(X_1, R^{s}\varphi_*(j_{!}F)) \Rightarrow H^{r+s}(X, j_{1!}F).
\]
按定理~\ref{thm:17.7}，$(R^{r}\varphi_*)(j_{!}F)$ 的茎可以在 $X/X_1$ 的几何纤维上计算.但 $x \in U$ 时 $\varphi$ 在 $\bar{x}$ 上的纤维由单点组成，而 $x \notin U$ 时 $j_{!}F$ 在该纤维上为零.由此推出 $r = 0$ 时 $R^{r}\varphi_*(j_{!}F) = j_{!}F$，否则为零.
\end{proof}

\begin{proposition}\label{prop:18.3}
(a) 层的短正合序列给出紧支集上同调群的长正合序列.

(b) 对 $U$ 的任何完备子簇 $Z$，有典范映射 $H^{r}_{Z}(U, F) \to H^{r}_{c}(U, F)$；对 $r = 0$，这些映射诱导同构
\[
H^{0}_{c}(U, F) \cong \varinjlim H^{0}_{Z}(U, F)
\]
($U$ 的完备子簇 $Z$ 上的极限).

(c) 若 $F$ 可构造，则 $H^{r}_{c}(U, F)$ 有限.
\end{proposition}

\begin{proof}
(a) 因为 $j_{!}$ 正合(\ref{prop:8.13})，$U$ 上层 的短正合序列
\[
0 \to F' \to F \to F'' \to 0
\]
给出 $X$ 上层 的短正合序列
\[
0 \to j_{!}F' \to j_{!}F \to j_{!}F'' \to 0,
\]
从而给出正合上同调序列
\[
\cdots \to H^{r}_{c}(U, F') \to H^{r}_{c}(U, F) \to H^{r}_{c}(U, F'') \to \cdots.
\]

(b) 设 $j \colon U \hookrightarrow X$ 是 $U$ 的完备化，$i \colon X \setminus U \hookrightarrow X$ 是互补的闭浸入.由正合序列(见命题~\ref{prop:8.15})
\[
0 \to j_{!}F \to j_{*}F \to i_{*}i^{*}j_{*}F \to 0,
\]
发现
\[
H^{0}_{c}(U, F) = \operatorname{Ker}(\Gamma(U, F) \to \Gamma(X \setminus U, i^{*}j_{*}F)).
\]
但由 $i^{*}$ 与 $j_{*}$ 的定义，看到\footnote{事实上，这不太正确，因为在形成逆像时需要层化(见 \S 8).}
\[
\Gamma(X \setminus U, i^{*}j_{*}F) = \varinjlim \Gamma(V \times_{X} U, F),
\]
其中极限遍历像含 $X \setminus U$ 的平展映射 $\varphi' \colon V \to X$.因此，$H^{0}_{c}(U, F)$ 是 $\Gamma(U, F)$ 中这样的截面子群：它在某个像含 $X \setminus U$ 的平展映射 $V \to X$ 对应的 $V \times_{X} U$ 上为零.

设 $Z$ 是 $U$ 的完备子簇，$s \in \Gamma_Z(U, F)$.于是 $s \in \Gamma(U, F)$ 且 $s|_{U \setminus Z} = 0$.因为 $Z$ 完备，它在 $X$ 中闭.现在 $s$ 在 $V \cap U$ 上为零，其中 $V$ 是 $Z$ 在 $X$ 中的(开)补集，这表明 $s \in H^{0}_{c}(U, F)$.我们已经证明
\[
\Gamma_Z(U, F) \subset H^{0}_{c}(U, F).
\]
反过来，设 $s \in H^{0}_{c}(U, F)$.则 $s$ 在某个像含 $X \setminus U$ 的平展映射 $\varphi' \colon V \to X$ 对应的 $V \times_{X} U$ 上为零.现在 $Z \overset{\mathrm{def}}{=} X \setminus \varphi'(V)$ 是 $U$ 的完备子簇，且 $U \setminus Z = U \cap \varphi'(V)$.因为 $V \times_{X} U \to U \cap \varphi'(V)$ 是平展覆盖且 $s|_{V \times_{X} U} = 0$，有 $s|_{U \setminus Z} = 0$.因此 $s \in \Gamma_Z(U, F)$，从而
\[
\bigcup_{Z} \Gamma_Z(U, F) = H^{0}_{c}(U, F).
\]
在上述证明过程中，我们证明了：对任何完备子簇 $Z \subset U$，$H^{0}_{Z}(U, F) \subset H^{0}_{c}(U, F)$.关于 $\delta$-函子的一般结果表明，态射 $H^{0}_{Z}(U, \cdot) \to H^{0}_{c}(U, \cdot)$ 唯一延拓为 $\delta$-函子的态射.显式地，假定该态射已经与连接同态相容地延拓到 $\leq r$ 维.给定 $F$，把它嵌入内射层 $F \hookrightarrow \mathcal{I}$，设 $F'$ 是商.有唯一的映射 $H^{r+1}_{Z}(U, F) \to H^{r+1}_{c}(U, F)$ 使下图交换：
\[
\begin{tikzcd}
\cdots \arrow[r] & H^{r}_{Z}(U, \mathcal{I}) \arrow[r] \arrow[d] & H^{r}_{Z}(U, F') \arrow[r] \arrow[d] & H^{r+1}_{Z}(U, F) \arrow[r] \arrow[d, dashed] & 0 \arrow[d] \\
\cdots \arrow[r] & H^{r}_{c}(U, \mathcal{I}) \arrow[r] & H^{r}_{c}(U, F') \arrow[r] & H^{r+1}_{c}(U, F) \arrow[r] & H^{r+1}_{c}(U, \mathcal{I})
\end{tikzcd}
\]
(c) 若 $F$ 可构造，则 $j_{!}F$ 也可构造，从而这由推论~\ref{cor:17.8} 推出.
\end{proof}

注意 $H^{r}_{c}(U, \cdot)$ 不是 $H^{0}_{c}(U, \cdot)$ 的第 $r$ 个右导出函子(从而不是导出函子).这令人遗憾：尽管 $H^{r}_{c}(U, \cdot)$ 不依赖于 $U$ 的完备化的选取，但似乎没有办法不先选取完备化就定义它：尽管 $H^{r}_{c}(U, \cdot)$ 是 $U$ 的内在对象，但似乎没有纯粹内在的定义.

\subsection*{具紧合支集的高阶直像}

利用 Nagata 定理的全部力量，可以定义具紧合支集的高阶直像(这里拓扑中的术语与逻辑所要求的术语一致).设 $\varphi \colon U \to S$ 是簇的正则映射(或概形的有限型分离态射)，$\bar{\varphi} \colon X \to S$ 是紧合态射，存在开浸入 $j \colon U \to X$，使得 $\varphi = \bar{\varphi} \circ j$ 且 $j(U)$ 在 $X$ 中稠密.对 $U$ 上的任何挠层 $F$，定义
\[
R^{r}\varphi_{!}F = R^{r}\bar{\varphi}_*(j_{!}F).
\]

\begin{proposition}\label{prop:18.4}
(a) 层 $R^{r}\varphi_{!}F$ 不依赖于分解 $\varphi = \bar{\varphi} \circ j$ 的选取.

(b) 层的短正合序列给出具紧合支集的高阶直像的长正合序列.

(c) 若 $F$ 可构造，则 $R^{r}\varphi_{!}F$ 也可构造.

(d) 对正则映射对 $U_1 \xrightarrow{\varphi_1} U_2 \xrightarrow{\varphi_2} S$ 与 $U_1$ 上的挠层 $F$，存在谱序列
\[
R^{r}\varphi_{2!}(R^{s}\varphi_{1!}F) \Rightarrow R^{r+s}(\varphi_2 \circ \varphi_1)_{!}(F).
\]
\end{proposition}

\begin{proof}
(a),(b),(c) 的证明与命题~\ref{prop:18.3} 中相应论断的证明类似.(d) 的证明因需要构造位于 $\varphi_1$ 的``完备化''之上的 $\varphi_2 \circ \varphi_1$ 的``完备化''而复杂---见 EC p.229.
\end{proof}

同样，$R^{r}\varphi_{!}$ 不是 $R^{0}\varphi_{!}$ 的第 $r$ 个右导出函子.更准确的记号应该是 $R^{r}_{c}\varphi_{*}$.
